\chapter{Introduction}
\setcounter{section}{0}

Sign Language Recognition (SLR) represents a highly interdisciplinary paradigm at the intersection
of digital accessibility, computer vision, and temporal sequence modelling, focused on capturing
fully-fledged, sophisticated natural languages where semantic meaning is conveyed simultaneously
rather than sequentially through distinct grammatical, syntactic, and visual-spatial structures.

\section{Problem description}

This project studies American Sign Language (ASL) recognition as an accessibility-oriented research
problem. Unlike static image classification, sign recognition is inherently temporal: a single
frame may show the handshape and approximate location of a sign, but it cannot represent the
direction of movement, the order of articulatory changes or the rhythm of execution. The sign must
therefore be interpreted from the ordered relationship between frames.

This temporal nature makes the problem substantially different from conventional object or gesture
recognition. Two signs may share a similar hand configuration in an individual frame, while
differing in trajectory, speed, orientation changes or relative hand position over time. In
addition, real videos introduce variability in signer appearance, camera viewpoint, lighting
conditions, signing style and execution speed. A recognition system must therefore learn not only
visual patterns, but also how those patterns evolve across a short sequence.

The task addressed in this work is isolated word recognition. The system receives a short video
clip and must assign it to a word gloss, that is, a written label used to annotate a sign. This
formulation is more limited than continuous sign language translation, where complete sentences,
grammar, co-articulation and broader linguistic context must be considered. However, isolated
recognition still captures a central challenge of SLR: identifying a meaningful visual unit from
motion and hand articulation.

\section{Research scope and approach}

The central research question of this work is whether compact visual representations can provide
enough information for temporal models to classify short ASL video clips into word glosses. The
project evaluates different neural architectures under reproducible experimental conditions, with
the objective of identifying a configuration that balances accuracy and computational cost. The
work does not aim to replace professional interpreters or to deliver a production-ready translator.

The experimental work is grounded on public ASL datasets, mainly WLASL \cite{li2020wlasl} and
MS-ASL \cite{joze2019msasl}. A mobile and server-based runtime prototype is also developed as a
proof of concept, although the main scientific contribution is the experimental evaluation of the
recognition approach rather than the deployment of an assistive product.

\section{Outline}

The remainder of the work is organised as follows:
\begin{itemize}
    \item Chapter 2 reviews the State of the Art, providing the theoretical foundation in computer
          vision and deep learning required for this work. It examines convolutional neural
          networks and pose estimation, specifically MediaPipe Hands, before addressing temporal
          modelling through recurrent networks, Long Short-Term Memory (LSTM) layers, Gated
          Recurrent Unit (GRU) layers, Transformers and attention mechanisms. The chapter also
          evaluates previous research in SLR to establish the context for the current approach.
    \item Chapter 3 defines the main objective, the specific objectives, and the validation methods
          used to determine whether each objective has been achieved.
    \item Chapter 4 outlines the methodology and project work plan. It details the iterative,
          experiment-driven approach adopted throughout development and delineates the main project
          phases. Finally, it addresses the technical resources and constraints that influenced the
          scope and design of the project.
    \item Chapter 5 presents the technical development and system architecture. It provides a
          rigorous account of the data preprocessing pipeline and the experimental framework
          established for training and evaluation. The chapter concludes with a description of a
          functional proof of concept, demonstrating a real-time implementation through a mobile
          app and a FastAPI-based inference server.
    \item Chapter 6 reports the experimental evaluation. It includes an architecture comparison and
          sequence-length analysis, followed by validation on the WLASL30 and MS-ASL50 datasets.
          The chapter also addresses cross-dataset validation, vocabulary scaling and the
          qualitative validation design, and concludes with a critical discussion of the main
          limitations observed during experimentation.
    \item Chapter 7 analyses the ethical implications and social impact of the project, structured
          around the Google Artificial Intelligence Principles and ethical framework. It provides a
          critical examination of accessibility, privacy and algorithmic bias, while addressing the
          requirements for responsible deployment and the importance of reproducibility. Finally,
          it evaluates the broader social value of assistive technologies through these ethical
          guidelines.
    \item Chapter 8 reflects on the academic journey at Universidad Francisco de Vitoria and
          relates the project to the competences developed during the degree.
    \item Chapter 9 presents the conclusions and future work, summarising the main findings, the
          limitations of the current prototype, and possible directions for extending the research.
\end{itemize}